\documentclass[a4paper]{article}
\usepackage[polish]{babel}
\usepackage[cp1250]{inputenc}
\usepackage[T1]{fontenc}
\usepackage[dvips]{graphicx}
\usepackage{colortbl}

\pagestyle{headings}
\textwidth      15.5cm
\oddsidemargin    .1cm
\evensidemargin   .1cm

\begin{document}

\thispagestyle{empty}

\begin{minipage}{5cm}
\includegraphics[scale=0.7]{salomon.eps}
\end{minipage}
\begin{minipage}{10cm}
\begin{center}
{\Large Salomon\\
System przetwarzania wiedzy\\}
\end{center}
\end{minipage}

\vspace*{0.5cm}

\hrulefill

\vspace*{1cm}


\begin{flushright}
{\Large \emph{Vision, ver 1.0} }
\end{flushright}

%\vspace{7cm}

\begin{flushleft}
{ \Large Historia wersji }\[\]
\begin{tabular}{|p{2.5cm}||p{1.5cm}|p{6cm}|p{3.5cm}|}
  \hline
  Data & Wersja & Opis & Autor \\
  \hline \hline
  25.04.2005 & 1.0 & utworzenie & \emph{Krzysztof Rajda}\\
  \hline
\end{tabular}
\end{flushleft}

%\[\] \[\] \[\] \[\] \[\] \[\] \[\] \[\] \[\]

% Zaczerpniete z MSF - osoba odpowiedzialna za dokument %
%\begin{flushleft}
%{ \Large Autor} \[\]
%\begin{tabular}{|p{3.5cm}||p{3cm}|p{4.5cm}|p{3cm}|}
%  \hline 
%  Leszek Grzanka & Tester,Programista & grzanka@student.agh.edu.pl & 31.03.2005 \\
%  \hline
%\end{tabular}
%\end{flushleft}

\newpage

\section{Wst�p}
Celem planu konstrukcji oprogramowania jest wyznaczenie r�l poszczeg�lnych cz�onk�w zespo�u oraz harmonogramu praca.
Przypisanie r�l pozwoli wyegzekwowa� brakuj�ce elementy systemu, od os�b za nie odpowiedzialnych oraz zlecenie poprawek, 
w wypadku, gdyby jaki� element jakie� elementy nie dzia�a�y zgodnie z dokumentacj�.
Stowrzenie harmonogramu pozwoli oceni� post�p prac i odpowiedzie� na pytanie, czy ko�cowy termin realizacji projektu jest osi�galny.

\section{Orgranizacja pojektu}

	\begin{flushleft}
		%{ \Large Historia wersji }\[\]
		\begin{tabular}{|l|l|}
  		\hline
  		Osoba & Rola wg. RUP \\
  		\hline \hline
  		Analyst & Krzysztof Rajda \\
  		\hline
  		Project Manager & Nikodem Jura\\
  		\hline
  		Process Engineer & Leszek Grzanka\\
  		\hline
  		Architect & Krzysztof Rajda\\
  		\hline
  		Developer & all \\
  		\hline
  		Tester & Jakub Ga�kowski \\
  		\hline
		\end{tabular}
	\end{flushleft}


\section{Terminarz projektu}
\subsection{Podzia� projektu na fazy i iteracje}
	\begin{flushleft}
		%{ \Large Historia wersji }\[\]
%		\begin{tabular}{|p{{0.75\textwidth}}|p{7cm}|p{7cm}|p{7cm}|p{7cm}|}
		\begin{tabular}{|l|l|p{3cm}|p{4cm}|p{3cm}|}
  		\hline
  		%\rowcolor[gray]{0.8}
  		Faza & Iteracja & G��wne cele & Produkty & Planowe rozpocz�cie/zako�czenie \\
  		\hline 
  		\hline
  		Inception &  I1 & Rozpoznanie problemu i otoczenia & Vision i SDP & 11.05.2005 \\
  		\hline
  		Eleboration & E1 & & Dokument SAD &  25.05.2005 \\
  		\hline
  		Construction & C1 & Stabilizacja Salomona & Stabilna wersja Salomona & 08.06.2005 \\
  		\hline
  		Construction & C2 & Poprawna obs�uga zbior�w danych & Wersja z pe�n� obs�uga zbior�w danych &  22.06.2005\\
  		\hline
  		Transition & T1 & Prezentacja dzia�ania zbior�w danych & Prezentacja Salomona &  29.06.2005\\
  		\hline
		\end{tabular}
	\end{flushleft}

\subsection{Diagram Gantt'a}
\section{Plan akceptacji produktu}
Akceptacja produktu b�dzie kilkuetapowa:
\begin{itemize}
	\item co nocne testy jednostkowe (doda� tutaj wszystkie rodzaje test�w?)
	\item przedstawienie klientowi do akceptacji produktu fazy \emph{elaboration},
	\item prezentacja klientowi ostatecznego produktu.
\end{itemize}
\end{document}

